% ===== PRÉAMBULE COMMUN — À COPIER TEL QUEL EN TÊTE DE CHAQUE FICHIER .tex =====
% Ne modifier QUE :
%   - les deux lignes \fancyhead[L] et \fancyhead[R]
%   - le bloc titre \begin{center}...\end{center} juste après \begin{document}
\documentclass[11pt,a4paper]{article}
\usepackage[utf8]{inputenc}
\usepackage[T1]{fontenc}
\usepackage{lmodern}
\usepackage[french]{babel}
\usepackage{amsmath,amssymb,mathtools}
\usepackage{icomma}
\usepackage{xcolor}
\usepackage{tikz}
\usepackage{pgfplots}
\usepackage[most]{tcolorbox}
\usepackage{enumitem}
\usepackage{array,booktabs,colortbl}
\usepackage[a4paper,margin=2.2cm,headheight=26pt,headsep=10pt]{geometry}
\usepackage{fancyhdr}
\usepackage[hidelinks]{hyperref}
\usetikzlibrary{arrows.meta,calc,angles,quotes,positioning,decorations.markings,intersections,backgrounds}
\pgfplotsset{compat=1.18}

\definecolor{primary}{RGB}{226,74,88}
\definecolor{primarydark}{RGB}{150,28,42}
\definecolor{methcol}{RGB}{40,90,150}
\definecolor{excol}{RGB}{0,135,90}
\definecolor{defcol}{RGB}{0,114,178}
\definecolor{attncol}{RGB}{200,40,55}
\definecolor{thmcol}{RGB}{0,140,120}
\definecolor{courbe}{RGB}{32,110,200}
\definecolor{courbeder}{RGB}{210,70,40}
\definecolor{tang}{RGB}{0,150,90}
\definecolor{grille}{RGB}{200,205,215}
\definecolor{plus}{RGB}{200,60,60}
\definecolor{moins}{RGB}{30,90,200}

\tcbset{boxbase/.style={enhanced,breakable,boxrule=0.9pt,arc=2.5pt,
  left=3mm,right=3mm,top=2.3mm,bottom=2.3mm,fonttitle=\bfseries,coltitle=white}}
\newtcolorbox{cle}[1][]{boxbase,colback=defcol!6,colframe=defcol,
  title={\textsc{À retenir}\ifx\relax#1\relax\relax\else\textmd{ : \itshape #1}\fi}}
\newtcolorbox{methode}[1][]{boxbase,colback=methcol!7,colframe=methcol,sharp corners,
  title={\textsc{Méthode}\ifx\relax#1\relax\relax\else\textmd{ --- \itshape #1}\fi}}
\newtcolorbox{exemplebox}[1][]{boxbase,colback=excol!5,colframe=excol,
  title={\textsc{Exemple}\ifx\relax#1\relax\relax\else\textmd{ : \itshape #1}\fi}}
\newtcolorbox{piege}[1][]{boxbase,colback=attncol!6,colframe=attncol,
  title={\textsc{Attention}\ifx\relax#1\relax\relax\else\textmd{ : \itshape #1}\fi}}
\newtcolorbox{exo}[1][]{boxbase,colback=primary!4,colframe=primary,
  title={\textsc{Exercice}\ifx\relax#1\relax\relax\else\textmd{~#1}\fi}}
\newtcolorbox{corr}[1][]{boxbase,colback=thmcol!5,colframe=thmcol,
  title={\textsc{Corrigé}\ifx\relax#1\relax\relax\else\textmd{~#1}\fi}}

\newcommand{\R}{\mathbb{R}}
\newcommand{\N}{\mathbb{N}}
\newcommand{\ds}{\displaystyle}
\newcommand{\tg}{\operatorname{tg}}
\newcommand{\cotg}{\operatorname{cotg}}
\newcommand{\arctg}{\operatorname{arctg}}
\newcommand{\arccotg}{\operatorname{arccotg}}
\newcommand{\limh}{\lim_{h\to 0}}

\pagestyle{fancy}\fancyhf{}
\fancyhead[L]{\small\textcolor{primarydark}{\textbf{Thème 5} — Tangente \& lecture graphique}}
\fancyhead[R]{\small\textcolor{primarydark}{\textsc{Tutoriel}}}
\fancyfoot[C]{\small\thepage}
\renewcommand{\headrulewidth}{0.8pt}

\begin{document}

\begin{center}
{\LARGE\bfseries\color{primarydark} Tangente à une courbe \& lecture graphique}\\[2pt]
{\color{primary}\rule{0.5\textwidth}{1.2pt}}\\[4pt]
{\small\itshape Fiche 05 — Dérivation \quad$\bullet$\quad Carnet d'entraînement \quad$\bullet$\quad Tutoriel}
\end{center}
\vspace{2mm}

\begin{cle}[l'équation de la tangente]
Soit $f$ dérivable en $x_0$. La \textbf{tangente} à la courbe de $f$ au point $P_0\big(x_0\,;f(x_0)\big)$ est la droite d'équation
\[
\boxed{\,y=(x-x_0)\,f'(x_0)+f(x_0)\,}
\]
\textbf{Provenance.} Une droite de pente $m$ passant par le point $(x_0;y_0)$ s'écrit $y=m(x-x_0)+y_0$. Pour la tangente, la pente vaut $m=f'(x_0)$ (c'est là tout le sens du nombre dérivé) et le point est $(x_0;f(x_0))$. On remplace : $y=f'(x_0)(x-x_0)+f(x_0)$.
\end{cle}

\begin{methode}[écrire l'équation de la tangente en $x_0$]
\begin{enumerate}[leftmargin=6mm,itemsep=1pt]
\item \textbf{Calculer $f(x_0)$} : l'ordonnée du point de contact.
\item \textbf{Dériver} : calculer $f'(x)$, puis évaluer $f'(x_0)$ (la pente).
\item \textbf{Remplacer} dans $y=(x-x_0)f'(x_0)+f(x_0)$.
\item \textbf{Développer} et réduire pour obtenir $y=mx+p$.
\end{enumerate}
\end{methode}

\begin{cle}[cas particuliers de pente]
\begin{itemize}[leftmargin=6mm,itemsep=1pt]
\item \textbf{Tangente horizontale} : elle a lieu là où $f'(x_0)=0$. Son équation est simplement $y=f(x_0)$ (souvent un extremum).
\item \textbf{Droites parallèles} : même pente, $m_1=m_2$.
\item \textbf{Droites perpendiculaires} : produit des pentes $m_1\cdot m_2=-1$, soit $m_2=-\dfrac{1}{m_1}$.
\item \textbf{Retrouver $f'(x_0)$} : si l'on connaît deux points $A$ et $B$ de la tangente, alors $f'(x_0)=\dfrac{y_B-y_A}{x_B-x_A}$ (pente entre deux points).
\end{itemize}
\end{cle}

\begin{center}
\begin{tikzpicture}[scale=1.0]
\begin{axis}[
  width=11.5cm,height=6.6cm,
  axis lines=middle,
  xlabel={$x$},ylabel={$y$},
  xlabel style={right},ylabel style={above},
  xmin=-2.4,xmax=3.2,ymin=-1.2,ymax=6.2,
  xtick={-2,-1,1,2,3},ytick={1,2,3,4,5},
  grid=both,grid style={grille!60,very thin},
  tick label style={font=\scriptsize},
]
% courbe f(x)=x^2
\addplot[courbe,very thick,domain=-2.3:2.6,samples=100]{x^2};
\addlegendentry{$\mathcal{C}_f:\;f(x)=x^2$}
% tangente en x0=1 : y=2x-1
\addplot[tang,very thick,domain=-1.0:3.0,samples=2]{2*x-1};
\addlegendentry{tangente en $P_0$}
% point P0
\addplot[only marks,mark=*,mark size=2pt,black]coordinates{(1,1)};
\node[anchor=south east,font=\small] at (axis cs:1,1){$P_0(1;1)$};
\end{axis}
\end{tikzpicture}
\end{center}

\begin{exemplebox}[tangente à $f(x)=x^2$ en $x_0=1$]
\textbf{1. Point :} $f(1)=1^2=1$, donc $P_0(1;1)$.\\
\textbf{2. Pente :} $f'(x)=2x$, donc $f'(1)=2$.\\
\textbf{3. Remplacer :} $y=(x-1)\cdot 2+1$.\\
\textbf{4. Développer :} $y=2x-2+1=\boxed{2x-1}$.
\end{exemplebox}

\begin{cle}[lecture graphique : relier $f$, $f'$ et $f''$]
\renewcommand{\arraystretch}{1.25}
\begin{center}
\begin{tabular}{@{}>{\raggedright\arraybackslash}p{6.2cm} c >{\raggedright\arraybackslash}p{6.2cm}@{}}
\toprule
\textbf{Sur le graphe de $f$} & & \textbf{Traduction} \\
\midrule
$f$ \textbf{monte} (croissante) & $\Longleftrightarrow$ & $f'>0$ \\
$f$ \textbf{descend} (décroissante) & $\Longleftrightarrow$ & $f'<0$ \\
$f$ présente un \textbf{extremum} (max ou min) & $\Longrightarrow$ & $f'=0$ (tangente horizontale) \\
$f$ est \textbf{convexe} (creux vers le haut, $\smallsmile$) & $\Longleftrightarrow$ & $f''>0$ \\
$f$ est \textbf{concave} (bosse vers le bas, $\smallfrown$) & $\Longleftrightarrow$ & $f''<0$ \\
\textbf{point d'inflexion} (changement de courbure) & $\Longrightarrow$ & $f''=0$ (et change de signe) \\
\bottomrule
\end{tabular}
\end{center}
\end{cle}

\end{document}
